\subsection{Markov Chains and Markov Processes}
% Transition kernels, semigroups, recurrence/transience, invariant distributions, ergodic behavior.

\begin{definition} \label{def:probability-vector-and-markov-matrix} \hlt{Probability Vectors and Markov Matrices}\\
A vector $\mathbf{x}\in\mathbb{R}^n$ is a \hlt{probability vector} if its entries are nonnegative and sum to $1$. Let $\mathbf{e}\defeq(1,\ldots,1)^t$. A matrix $A=[a_{ij}]\in\operatorname{Mat}(n,n,\mathbb{R})$ is a \hlt{Markov matrix} if $a_{ij}\geq0$ for all $i,j$ and $A\mathbf{e}=\mathbf{e}$; equivalently, every row sum is $1$.
\end{definition}

\begin{lemma} \label{lem:markov-matrix-basic-properties} \hlt{Basic Properties of Markov Matrices}\\
If $A\in\operatorname{Mat}(n,n,\mathbb{R})$ is Markov, then $1\in\operatorname{spec}(A)$, and each $A^p$ is Markov with entries in $[0,1]$.
\begin{proof}
Since $A\mathbf{e}=\mathbf{e}$, $1$ is an eigenvalue. Inductively, $A^p$ has nonnegative entries and $A^p\mathbf{e}=\mathbf{e}$, so $A^p$ is Markov. Each nonnegative row sums to $1$, so all entries lie in $[0,1]$.\\
\hlt{Structural Commentary}: Markov structure is closed under powers.\\
\hlt{Key Insight}: Nonnegative row sums of $1$ bound every entry.
\end{proof}
\end{lemma}

\begin{theorem} \label{thm:markov-matrix-unit-circle-semisimple} \hlt{Unit-Circle Spectrum of Markov Matrices}\\
If $A\in\operatorname{Mat}(n,n,\mathbb{R})$ is Markov, then $\rho(A)=1$, and $\dim E_{\lambda}(A)=m_A(\lambda)$ for every $\lambda\in\operatorname{spec}(A)$ with $\abs{\lambda}=1$.
\begin{proof}
Lemma \ref{lem:markov-matrix-basic-properties} makes $A$ power bounded. Theorem \ref{thm:matrix-power-bounded-jordan-criterion} gives $\rho(A)\leq1$ and $\dim E_{\lambda}(A)=m_A(\lambda)$ on $\abs{\lambda}=1$. Since $A\mathbf{e}=\mathbf{e}$, $1\in\operatorname{spec}(A)$, so $\rho(A)\geq1$. Thus $\rho(A)=1$.\\
\hlt{Structural Commentary}: Bounded Markov powers rule out unit-circle Jordan growth.\\
\hlt{Key Insight}: $A\mathbf{e}=\mathbf{e}$ forces the eigenvalue $1$.
\end{proof}
\end{theorem}

\begin{theorem} \label{thm:positive-markov-stationary-distribution} \hlt{Positive Markov Matrices Converge to a Stationary Distribution}\\
Let $n\geq2$ and let $A\in\operatorname{Mat}(n,n,\mathbb{R})$ be a Markov matrix with positive entries. Then there is a unique probability vector $\mathbf{x}$ with $A^t\mathbf{x}=\mathbf{x}$; its entries are positive and $\lim\limits_{p\to\infty}A^p=\mathbf{e}\mathbf{x}^t$.
\begin{proof}
Let $A\mathbf{y}=\lambda\mathbf{y}$ with $\abs{\lambda}=1$, and choose $k$ with $\abs{y_k}=\norm{\mathbf{y}}_{\infty}$.\\
Positivity and the $k$th row sum give $\norm{\mathbf{y}}_{\infty} =\abs{\lambda y_k} =\abs{\sum\limits_{j=1}^{n}a_{kj}y_j} \leq\sum\limits_{j=1}^{n}a_{kj}\abs{y_j} \leq\norm{\mathbf{y}}_{\infty}$.\\
Because every $a_{kj}$ is positive, equality in the last inequality forces $\abs{y_j}=\norm{\mathbf{y}}_{\infty}$ for all $j$. Equality in the triangle inequality then forces all nonzero $y_j$ to have the same argument. Hence $\mathbf{y}=c\mathbf{e}$, and $A\mathbf{e}=\mathbf{e}$ gives $\lambda=1$. Thus $E_1(A)=\operatorname{span}(\mathbf{e})$ and no other eigenvalue lies on the unit circle.\\
By Theorem \ref{thm:markov-matrix-unit-circle-semisimple}, the Jordan form is $[1]\oplus J$ with $\rho(J)<1$. Choose $S=[\mathbf{e}\ Y]$ such that $A=S([1]\oplus J)S^{-1}$, and let $\mathbf{x}^t$ be the first row of $S^{-1}$. Then $\mathbf{x}^tA=\mathbf{x}^t$ and $\mathbf{x}^t\mathbf{e}=1$. Since $J^p\to0$, $A^p=S([1]\oplus J^p)S^{-1}\to \mathbf{e}\mathbf{x}^t$.\\
The nonnegative powers $A^p$ imply $\mathbf{e}\mathbf{x}^t\geq0$, so $\mathbf{x}\geq0$. Since $A^t\mathbf{x}=\mathbf{x}$, $\mathbf{x}\neq0$, and $A^t$ has positive entries, in fact $\mathbf{x}>0$. Finally, $\dim E_1(A^t)=\dim E_1(A)=1$ because $A-I_n$ and $(A-I_n)^t$ have the same rank, so $\mathbf{x}^t\mathbf{e}=1$ gives uniqueness among probability vectors.\\
\hlt{Structural Commentary}: Positivity leaves one simple unit-circle block; all other blocks decay.\\
\hlt{Key Insight}: The limit is the rank-one projection $\mathbf{e}\mathbf{x}^t$.
\end{proof}
\end{theorem}

\begin{figure}[H]
  \centering
  \begin{tikzpicture}[
  font=\small,
  block/.style={draw=blue!65!black, fill=blue!7, rounded corners=2pt, minimum width=2.25cm, minimum height=0.64cm, align=center},
  decay/.style={draw=green!45!black, fill=green!8, rounded corners=2pt, minimum width=2.25cm, minimum height=0.64cm, align=center},
  arrow/.style={-{Stealth[length=2mm]}, thick, gray!75}
]
  \node[block] (A) at (0,0) {$A=S([1]\oplus J)S^{-1}$};
  \node[block] (line) at (3.6,0.7) {$\operatorname{span}(\mathbf{e})$};
  \node[decay] (rest) at (3.6,-0.7) {$\rho(J)<1$};
  \node[block] (limit) at (7.2,0) {$A^p\to\mathbf{e}\mathbf{x}^t$};

  \draw[arrow] (A.east) -- (line.west);
  \draw[arrow] (A.east) -- (rest.west);
  \draw[arrow] (line.east) -- (limit.west);
  \draw[arrow] (rest.east) -- (limit.west);
\end{tikzpicture}

  \caption{The line $\operatorname{span}(\mathbf{e})$ remains; all other Jordan directions decay, so $A^p\to\mathbf{e}\mathbf{x}^t$.}
  \label{fig:positive-markov-stationary-limit}
\end{figure}

The preceding results are matrix facts. In probability language, a row-stochastic matrix becomes a transition matrix for a discrete-time chain.

\begin{definition} \label{def:homogeneous-markov-chain-transition-matrix} \hlt{Markov Chain and Transition Matrix}
\begin{enumerate}[label=(\roman*)]
\setlength{\itemsep}{0pt}
\item Let $\mathcal{X}$ be countable and let $(X_n)_{n\geq0}$ take values in $\mathcal{X}$.
\item Process is \hlt{Markov chain} if $\mathbb{P}(X_n=x\mid X_0,\ldots,X_{n-1})=\mathbb{P}(X_n=x\mid X_{n-1})$ $\forall$ admissible $n$ and $x\in\mathcal{X}$.
\item It is \hlt{homogeneous} if $p_{ij}\defeq\mathbb{P}(X_{n+1}=j\mid X_n=i)$ is independent of $n$.
\item The matrix $P=[p_{ij}]$ is the \hlt{transition matrix}; $p_{ij}\geq0$ and $\sum_{j\in\mathcal{X}}p_{ij}=1$ for each $i$.
\end{enumerate}
\end{definition}

\begin{figure}[H]
  \centering
  \begin{tikzpicture}[>=Stealth,node distance=1.1cm,every node/.style={font=\small}]
  \node (x1) {$X_1$};
  \node[right=of x1] (x2) {$X_2$};
  \node[right=of x2] (x3) {$X_3$};
  \node[right=of x3] (dots) {$\cdots$};
  \node[right=of dots] (xn) {$X_n$};
  \node[right=of xn] (dots2) {$\cdots$};
  \draw[->,thick] (x1) -- (x2);
  \draw[->,thick] (x2) -- (x3);
  \draw[->,thick] (x3) -- (dots);
  \draw[->,thick] (dots) -- (xn);
  \draw[->,thick] (xn) -- (dots2);
\end{tikzpicture}

  \caption{A Markov chain remembers the current state, not the full past.}
  \label{fig:markov-chain-dag}
\end{figure}

\begin{example} \label{ex:random-walk-absorbing-barriers} \hlt{Random Walk with Absorbing Barriers}\\
Let $\mathcal{X}=\{1,\ldots,N\}$ and $q=1-p$.
The absorbing random walk has $p_{11}=p_{NN}=1$, $p_{i,i+1}=p$ and $p_{i,i-1}=q$ for $2\leq i\leq N-1$, and all other transition probabilities equal to $0$.
\end{example}

\begin{definition} \label{def:n-step-transition-probabilities} \hlt{$n$-Step Transition Probabilities}\\
For a homogeneous Markov chain, the \hlt{$n$-step transition probability} from $i$ to $j$ is\\
$p_{ij}(n)\defeq\mathbb{P}(X_{m+n}=j\mid X_m=i)$; set $P_n=[p_{ij}(n)]$.
\end{definition}

\begin{theorem} \label{thm:chapman-kolmogorov-discrete} \hlt{Chapman-Kolmogorov Equations}\\
For a homogeneous Markov chain on countable $\mathcal{X}$, for all $i,j\in\mathcal{X}$ and $m,n\geq0$, $p_{ij}(m+n)=\sum_{k\in\mathcal{X}} p_{ik}(m)p_{kj}(n)$.\\
Equivalently, $P_{m+n}=P_mP_n$ for $P_n=P^n$.
\begin{proof} By the law of total probability and the Markov property,\\
$\mathbb{P}(X_{m+n}=j\mid X_0=i)=\sum_{k\in\mathcal{X}}\mathbb{P}(X_{m+n}=j\mid X_m=k)\mathbb{P}(X_m=k\mid X_0=i)$.\\
Homogeneity identifies the two factors as $p_{kj}(n)$ and $p_{ik}(m)$, giving the scalar identity.\\
Matrix multiplication gives $P_{m+n}=P_mP_n$, and $P_1=P$ yields $P_n=P^n$ by induction.\\
\hlt{Structural Commentary}: Transition over $m+n$ steps factors through the intermediate state after $m$ steps.\\
\hlt{Key Insight}: Conditioning over intermediate states is matrix multiplication.
\end{proof}
\end{theorem}

\begin{lemma} \label{lem:markov-marginal-evolution} \hlt{Marginal Evolution}\\
Let $\mu_n$ be the row distribution of $X_n$, so $\mu_n(j)=\mathbb{P}(X_n=j)$.\\
For a homogeneous chain with transition matrix $P$, $\mu_n=\mu_0P^n$.
\begin{proof}
For each $j$, $\mu_n(j)=\sum_{i\in\mathcal{X}}\mathbb{P}(X_n=j\mid X_0=i)\mathbb{P}(X_0=i)=\sum_{i\in\mathcal{X}}\mu_0(i)p_{ij}(n)$.\\
Since $P_n=P^n$, this is the $j$th component of $\mu_0P^n$.\\
\hlt{Structural Commentary}: Marginals evolve by multiplying the initial law by the transition semigroup.\\
\hlt{Key Insight}: Conditioning on $X_0$ turns distribution evolution into row-vector matrix multiplication.
\end{proof}
\end{lemma}

\begin{definition} \label{def:markov-state-communication} \hlt{Communication Classes}
\begin{enumerate}[label=(\roman*)]
\setlength{\itemsep}{0pt}
\item State $i$ \hlt{reaches} $j$, written $i\to j$, if $p_{ij}(n)>0$ for some $n\geq0$.
\item States $i,j$ \hlt{communicate}, written $i\leftrightarrow j$, if $i\to j$ and $j\to i$.
\item A \hlt{communication class} is an equivalence class under $\leftrightarrow$.
\item The chain is \hlt{irreducible} if all states lie in one communication class.
\item A set $C$ is \hlt{closed} if $p_{ij}=0$ for all $i\in C$ and $j\notin C$; a singleton closed class is an \hlt{absorbing state}.
\end{enumerate}
\end{definition}

\begin{proposition} \label{prop:communication-equivalence-classes} \hlt{Communication Classes Partition the State Space}\\
The relation $\leftrightarrow$ on $\mathcal{X}$ is an equivalence relation. Hence $\mathcal{X}=\bigsqcup_{\alpha\in A}C_\alpha$, each $C_\alpha$ is a communication class.
\begin{proof} Reflexivity follows from $p_{ii}(0)=1$, and symmetry is part of the definition of $\leftrightarrow$.\\
For transitivity, if $i\leftrightarrow j$, $j\leftrightarrow k$, choose $m,n,r,s\geq0$ with $p_{ij}(m)>0$, $p_{jk}(n)>0$, $p_{kj}(r)>0$, $p_{ji}(s)>0$.\\
Then $p_{ik}(m+n)\geq p_{ij}(m)p_{jk}(n)>0$ and $p_{ki}(r+s)\geq p_{kj}(r)p_{ji}(s)>0$, so $i\leftrightarrow k$.\\
Equivalence classes of $\leftrightarrow$ are disjoint and cover $\mathcal{X}$.\\
\hlt{Structural Commentary}: Communication classes are the equivalence classes generated by mutual reachability.\\
\hlt{Key Insight}: Positive transition paths concatenate by Chapman-Kolmogorov.
\end{proof}
\end{proposition}

\begin{definition} \label{def:recurrent-transient-state} \hlt{Recurrent and Transient States}
\begin{enumerate}[label=(\roman*)]
\setlength{\itemsep}{0pt}
\item Let $\tau_i^+\defeq\inf\{n\geq1:X_n=i\}$ be the first return time to $i$.
\item State $i$ is \hlt{recurrent} or \hlt{persistent} if $\mathbb{P}_i(\tau_i^+<\infty)=1$.
\item State $i$ is \hlt{transient} if $\mathbb{P}_i(\tau_i^+<\infty)<1$.
\end{enumerate}
\end{definition}

\begin{theorem} \label{thm:recurrence-transition-sum-criterion} \hlt{Recurrence Criterion}
\begin{enumerate}[label=(\roman*)]
\setlength{\itemsep}{0pt}
\item State $i$ is \hlt{recurrent} iff $S_i\defeq\sum_{n=0}^{\infty}p_{ii}(n)=\infty$.
\item State $i$ is \hlt{transient} iff $S_i<\infty$.
\end{enumerate}
\begin{proof} Let $N_i\defeq\sum_{n=0}^{\infty}\mathbf{1}\{X_n=i\}$ and $r_i\defeq\mathbb{P}_i(\tau_i^+<\infty)$. By Tonelli, $\mathbb{E}_iN_i=\sum_{n=0}^{\infty}p_{ii}(n)=S_i$.\\
By the Markov property after returns to $i$, $\mathbb{P}_i(N_i\geq k)=r_i^{k-1}$ for $k\geq1$, hence $\mathbb{E}_iN_i=\sum_{k=1}^{\infty}r_i^{k-1}$.
\begin{enumerate}[label=(\roman*)]
\setlength{\itemsep}{0pt}
\item[$(i)$] ($\Rightarrow$) If $i$ is recurrent, then $r_i=1$, so $S_i=\mathbb{E}_iN_i=\infty$.
\item[$(i)$] ($\Leftarrow$) If $S_i=\infty$, then $r_i<1$ would give $\mathbb{E}_iN_i=(1-r_i)^{-1}<\infty$, a contradiction; hence $r_i=1$.
\item[$(ii)$] ($\Rightarrow$) If $i$ is transient, then $r_i<1$, so $S_i=\mathbb{E}_iN_i=(1-r_i)^{-1}<\infty$.
\item[$(ii)$] ($\Leftarrow$) If $S_i<\infty$, then $r_i=1$ would give $\mathbb{E}_iN_i=\infty$, a contradiction; hence $r_i<1$.
\end{enumerate}
\hlt{Structural Commentary}: Recurrence is detected by the expected total number of returns.\\
\hlt{Key Insight}: Visits to $i$ repeat with geometric return probability $r_i$.
\end{proof}
\end{theorem}

\begin{definition} \label{def:markov-period-ergodic-state} \hlt{Period and Ergodicity}
\begin{enumerate}[label=(\roman*)]
\setlength{\itemsep}{0pt}
\item For recurrent $i$, it is \hlt{positive recurrent} if $\mathbb{E}_i\tau_i^+<\infty$ and \hlt{null recurrent} otherwise.
\item The \hlt{period} of $i$ is $d(i)\defeq\gcd\{n\geq1:p_{ii}(n)>0\}$.
\item State $i$ is \hlt{aperiodic} if $d(i)=1$.
\item State $i$ is \hlt{ergodic} if it is recurrent, positive recurrent, and aperiodic.
\item The chain is \hlt{ergodic} if every state is ergodic.
\end{enumerate}
\end{definition}

\begin{definition} \label{def:markov-stationary-limiting-distributions} \hlt{Stationary and Limiting Distributions}
\begin{enumerate}[label=(\roman*)]
\setlength{\itemsep}{0pt}
\item A probability row vector $\pi$ is \hlt{stationary} or \hlt{invariant} if $\pi=\pi P$.
\item A finite chain has \hlt{limiting distribution} $\pi$ if $p_{ij}(n)\to\pi_j$ for all states $i,j$.
\end{enumerate}
Equivalently, every row of $P^n$ converges to $\pi$.
\end{definition}

\begin{theorem} \label{thm:finite-ergodic-chain-convergence} \hlt{Finite Ergodic Chain Convergence}\\
Let $P$ be the transition matrix of a finite irreducible aperiodic Markov chain on $\mathcal{X}$.\\
There is a unique $\pi\in\Delta(\mathcal{X})$ with $\pi P=\pi$, and $P^n\to\mathbf{e}\pi$.\\
For every bounded $g:\mathcal{X}\to\mathbb{R}$, $N^{-1}\sum_{n=1}^{N}g(X_n)\to\pi(g)\defeq\sum_{j\in\mathcal{X}}g(j)\pi_j$ almost surely.
\begin{proof} Finite irreducibility and aperiodicity imply $P$ is primitive: $P^m>0$ for some $m\geq1$.\\
Row-vector form of Theorem \ref{thm:positive-markov-stationary-distribution} on $A=P^m$, gives unique $\pi\in\Delta(\mathcal{X})$ with $\pi P^m=\pi$, $(P^m)^k\to\mathbf{e}\pi$.\\
Since $(\pi P)P^m=\pi P$ and $\pi P\in\Delta(\mathcal{X})$, uniqueness gives $\pi P=\pi$.\\
For $n=mk+r$ with $0\leq r<m$, $P^n=(P^m)^kP^r\to\mathbf{e}\pi P^r=\mathbf{e}\pi$.\\
The finite-state Markov-chain strong law yields $N^{-1}\sum_{n=1}^{N}g(X_n)\to\pi(g)$ almost surely for bounded $g$.\\
\hlt{Structural Commentary}: Irreducibility gives one communicating class; aperiodicity removes cyclic obstruction.\\
\hlt{Key Insight}: Primitivity makes $P^n$ converge to the rank-one projection $\mathbf{e}\pi$.
\end{proof}
\end{theorem}

\begin{definition} \label{def:detailed-balance} \hlt{Detailed Balance}\\
Let $P=[p_{ij}]$ be a transition matrix on $\mathcal{X}$.\\
A probability row vector $\pi$ satisfies \hlt{detailed balance} for $P$ if $\pi_i p_{ij}=\pi_jp_{ji}$ for all $i,j\in\mathcal{X}$.
\end{definition}

\begin{theorem} \label{thm:detailed-balance-stationary} \hlt{Detailed Balance Implies Stationarity}\\
Let $P=[p_{ij}]$ be row-stochastic on countable $\mathcal{X}$, and let $\pi\in\Delta(\mathcal{X})$. If $\pi_i p_{ij}=\pi_jp_{ji}$ $\forall i,j\in\mathcal{X}$, then $\pi P=\pi$.
\begin{proof} $\forall j\in\mathcal{X}$, $(\pi P)_j=\sum_{i\in\mathcal{X}}\pi_i p_{ij}=\sum_{i\in\mathcal{X}}\pi_jp_{ji}=\pi_j\sum_{i\in\mathcal{X}}p_{ji}=\pi_j$, since row $j$ of $P$ sums to $1$. QED.\\
\hlt{Structural Commentary}: Detailed balance gives pairwise cancellation of probability fluxes.\\
\hlt{Key Insight}: Reversibility is stronger than stationarity.
\end{proof}
\end{theorem}

\begin{remark} \label{rem:stationary-does-not-imply-convergent} \hlt{Stationary Need Not Mean Convergent}\\
For the deterministic three-cycle on $\{1,2,3\}$, set $p_{12}=p_{23}=p_{31}=1$ and all other $p_{ij}=0$.\\
Then $\pi=(1/3,1/3,1/3)$ satisfies $\pi P=\pi$, but $P^{n+3}=P^n$, so $P^n$ does not converge.
\end{remark}

\begin{method} \label{method:markov-transition-mle} \hlt{Estimating a Finite Transition Matrix}
\begin{enumerate}[label=(\roman*)]
\setlength{\itemsep}{0pt}
\item Observe finite homogeneous chain $X_0,\ldots,X_n$, let $n_{ij}\defeq\sum_{t=0}^{n-1}\mathbf{1}\{X_t=i,X_{t+1}=j\}$, $n_i\defeq\sum_j n_{ij}$.
\item The row likelihood is proportional to $\prod_j p_{ij}^{n_{ij}}$ subject to $\sum_jp_{ij}=1$.
\item If $n_i>0$, the rowwise MLE is $\widehat{p}_{ij}=n_{ij}/n_i$.
\item If $n_i=0$, row $i$ is not identifiable without a convention or prior.
\end{enumerate}
\end{method}

\begin{proposition} \label{prop:two-state-chain-limit} \hlt{Two-State Chain Limit}\\
Let $0<a,b<1$ and $P=\begin{bmatrix}
1-a&a\\
b&1-b
\end{bmatrix}$.\\
Set $\pi=\left(b/(a+b),a/(a+b)\right)$. Then $\pi P=\pi$, $\pi$ is unique, and $P^n\to\mathbf{e}\pi=
\begin{bmatrix}
\frac{b}{a+b}&\frac{a}{a+b}\\
\frac{b}{a+b}&\frac{a}{a+b}
\end{bmatrix}$.
\begin{proof} Solving $\pi P=\pi$ with $\pi\mathbf{e}=1$ gives $\pi=\left(b/(a+b),a/(a+b)\right)$.\\
The characteristic polynomial is $(\lambda-1)(\lambda-(1-a-b))$, so $\operatorname{spec}(P)=\{1,1-a-b\}$ and $\abs{1-a-b}<1$.\\
Thus stationary component is projected onto $\mathbf{e}\pi$, transverse component decays geometrically, giving $P^n\to\mathbf{e}\pi$.\\
\hlt{Structural Commentary}: The dynamics split into the stationary line and one transverse eigendirection.\\
\hlt{Key Insight}: The non-unit eigenvalue $1-a-b$ controls convergence.
\end{proof}
\end{proposition}
